\documentclass[12pt]{article}
\usepackage[utf8]{inputenc}
\usepackage[T1]{fontenc}
\usepackage{hyperref}
\title{Cuneiform Luwian Language Resources for K-12 Curriculum Development}
\author{Compiled by Claude Sonnet 5.5 (claude-sonnet-5-5) for Andr\'as Kornai}
\date{}
\begin{document}
\maketitle

\section{Introduction}
Cuneiform Luwian (ISO 639-3 xlu; Glottocode cune1239) is an extinct Anatolian Indo-European language. It is attested in cuneiform texts from the Hittite archives, which are mostly ritual texts. It has no native speakers. The resources found are scholarly, and there are no school materials. A K-12 use would be a heritage and ancient-languages enrichment unit, not a language-learning curriculum.

\section{Primary Educational Resources}
I found no primers, textbooks or school programmes for Cuneiform Luwian.

\section{Linguistic and Reference Materials}
\begin{itemize}
\item Yakubovich and Rieken, \emph{A Grammar of the Luwian Language} (Harrassowitz, StBoT 75). It is a print work and the publisher lists it as released in February 2026. It covers both the cuneiform and hieroglyphic corpora. \url{https://www.harrassowitz-verlag.de/A_Grammar_of_the_Luwian_Language/title_9153.ahtml}
\item H. Craig Melchert, \emph{A Dictionary of Cuneiform Luvian} (print), listed by a bookseller: \url{https://www.ebook.de/de/product/53495205/h_craig_melchert_a_dictionary_of_cuneiform_luvian.html}
\item Cuneiform Luvian Lexicon, an online provisional index of attested lexemes (University of Gothenburg): \url{https://databaser.ub.gu.se/cuneiform-luvian-lexicon/188893}
\item eDiAna online corpus with interlinear glosses (LMU Munich). Example texts: \url{https://www.ediana-remake.gwi.uni-muenchen.de/corpus/cuneiform/1069}, \url{https://www.ediana-remake.gwi.uni-muenchen.de/corpus/cuneiform/1166}, \url{https://www.ediana-remake.gwi.uni-muenchen.de/corpus/cuneiform/1203}
\item eDiAna dictionary entry example: \url{https://www.ediana-remake.gwi.uni-muenchen.de/dictionary/1527}
\end{itemize}
The standard text edition is Starke's \emph{Die keilschrift-luwischen Texte in Umschrift} (StBoT 30, 1985). The search results cited it but did not link to it.

\section{Religious and Community Resources}
The surviving texts are mainly Hittite-archive ritual texts. I found no community organisations and no modern religious resources.

\section{Digital and Audio Resources}
The eDiAna corpus and the Gothenburg lexicon are the digital resources found. I found no audio or video resources, and none are expected because no pronunciation tradition survives.

\section{Curriculum Development Context}
Realistic classroom uses are a unit on Bronze Age Anatolia and writing systems, reading a few glossed lines from eDiAna as a decipherment exercise, and comparison with Hittite and Hieroglyphic Luwian. Teachers would need to write all of the material themselves.
\end{document}